% TOP_PROBLEM: {"id":30007675,"problem_number":"LOCAL-30007675","title":"Merger effects beyond prices","category_id":21,"category":{"id":21,"name":"economics","display_name":"Economics","description":"Open economic theory statements, published research agendas, and proposed scoped specifications; question provenance and historical priority are qualified per record.","slug":"economics","order_index":21,"created_at":"2026-10-08T00:00:00Z"},"status":"research_question","external_url":null,"legacy_ids":[],"published":false,"statement":"Estimate merger effects on independently audited quality and available product variety over one, three and five years, including entry and discontinuation.","economics":{"original_id":164,"field":"Industrial organization and competition","kind":"empirical","formalization_basis":"adapted_research_question","source_status":"research_agenda","priority_status":"earliest_found","priority_note":"Earliest relevant explicit question or agenda passage located in this audit. No claim of first-ever formulation; earlier versions and later solutions were not exhaustively traced.","earliest_verified_reference":{"authors":["Carl Shapiro","Ali Yurukoglu"],"title":"Trends in Competition in the United States: What Does the Evidence Show?","year":2024,"url":"https://web.stanford.edu/~ayurukog/trendsincompetition.pdf","doi":"10.3386/w32762","locator":"Section 6, printed pp. 35–36, merger-retrospective paragraphs","evidence_summary":"Calls for additional evidence on merger effects beyond prices, including innovation."},"current_reference":{"authors":["Carl Shapiro","Ali Yurukoglu"],"title":"Trends in Competition in the United States: What Does the Evidence Show?","year":2024,"url":"https://web.stanford.edu/~ayurukog/trendsincompetition.pdf","doi":"10.3386/w32762","locator":"Section 6, printed pp. 35–36, merger-retrospective paragraphs","evidence_summary":"Calls for additional evidence on merger effects beyond prices, including innovation."},"formalization_note":"The source verifies a broader question or research agenda; the population, estimand, equations and restrictions here are our scoped adaptation. This is not a quoted canonical conjecture, and the audit does not prove contemporary unsolved status for every special case."},"statusQualification":"Published question or research agenda verified at the cited locator. The mathematical specification is adapted for this catalogue and includes additional assumptions; source publication does not establish first-ever priority or certify this exact model remains unsolved. The source verifies a broader question or research agenda; the population, estimand, equations and restrictions here are our scoped adaptation. This is not a quoted canonical conjecture, and the audit does not prove contemporary unsolved status for every special case.","statusReviewedAt":"2026-10-08"}
% FIRST_POSTED: 2026-10-08
% ENTRY_KIND: statement_only
% !TeX program = lualatex
\documentclass[11pt]{article}
\usepackage{fontspec}
\usepackage[a4paper,margin=25mm]{geometry}
\usepackage{amsmath,amssymb,mathtools}
\usepackage{xurl}
\usepackage[hidelinks]{hyperref}
\newcommand{\sref}[2]{\hyperlink{source:#1:#2}{[S#2]}}
\setlength{\emergencystretch}{3em}
\begin{document}
% problemId:30007675
\section*{Merger effects beyond prices}\label{30007675}
\subsection{Definitions and mathematical statement}
Fix the population of horizontal mergers in United States grocery-product markets closing during a specified five-year cohort. For merger \(m\), define two regimes \(a\in\{0,1\}\), where \(a=1\) implements the merger and \(a=0\) keeps the firms independent. Let \(Q_m(a,h)\) be an independently audited quality index and \(V_m(a,h)\) the number of distinct product variants available in its premerger geographic market at horizon \(h\) years. Normalize quality on a common scale and fix the variant definition before seeing outcomes. The primary targets are \(\tau_Q(h)=E[Q_m(1,h)-Q_m(0,h)]\) and \(\tau_V(h)=E[V_m(1,h)-V_m(0,h)]\), for \(h\in\{1,3,5\}\); expectation averages over the defined merger cohort. Include entry, exit and product redesign within each regime, instead of holding the original product set fixed. Identify these effects using comparable markets under an explicit conditional counterfactual-trends assumption, or report sensitivity bounds when that assumption is relaxed. Quality sampling must not depend only on surviving products, since discontinuation is itself an outcome. Quantify whether quality and variety responses differ with premerger substitution and cost complementarities. This specifies a longitudinal empirical agenda: it does not assume that every merger lowers quality, and it distinguishes changes in product count from changes in consumer-valued variety.

\par\textbf{Formulation and scope.} The source verifies a broader question or research agenda; the population, estimand, equations and restrictions here are our scoped adaptation. This is not a quoted canonical conjecture, and the audit does not prove contemporary unsolved status for every special case.

\par\textbf{Open-question provenance.} An explicit question or research agenda was verified at the source locator. This does not by itself certify that every later version remains unresolved \sref{30007675}{1}.

\par\textbf{Historical priority.} Earliest relevant explicit question or agenda passage located in this audit. No claim of first-ever formulation; earlier versions and later solutions were not exhaustively traced.

\subsection{Short English statement}
Estimate merger effects on independently audited quality and available product variety over one, three and five years, including entry and discontinuation.

\subsection{Sources}
\begin{itemize}\raggedright
\item[\textnormal{[S1]}]\hypertarget{source:30007675:1}{} Carl Shapiro, Ali Yurukoglu. 2024. Trends in Competition in the United States: What Does the Evidence Show?. Section 6, printed pp. 35–36, merger-retrospective paragraphs.
\url{https://web.stanford.edu/~ayurukog/trendsincompetition.pdf}.
DOI: 10.3386/w32762.
{\small\textit{Earliest verified question/agenda reference; historical priority qualified below. Calls for additional evidence on merger effects beyond prices, including innovation.}}
\end{itemize}
\end{document}
